\documentclass[output=paper,colorlinks,citecolor=brown
% ,hidelinks
% showindex
]{langscibook}
\ChapterDOI{10.5281/zenodo.20285210}
\author{Ronald P. Schaefer\orcid{}\affiliation{Southern Illinois University Edwardsville} and        Francis Egbokhare\orcid{}\affiliation{University of Ibadan}}

\title{Beyond the transitive prototype: Bodily-condition predications in Emai}
\abstract{We assess transitivity prominence from the perspective of its periphery for an Edoid language of Nigeria. Emai exhibits a set of simple verb predications that assume a two-position transitive shape and express a bodily condition or an abstract condition. They permit a human participant in one position and a body-part noun or abstract noun in the other. They deviate from the transitive prototype, whose two participants exhibit meanings that are maximally distinct: a volitional instigator vs. a non-volitional entity showing material effect. Deviation is evident in the failure of participant distinctness and failure of material effect. Only the general condition on two-participant form remains of the transitive prototype. We also consider whether bodily-condition predications might belong to a grammatical domain other than transitivity. Copulas in Emai are considered since English translation of most examples employ a form of BE. There are five copula constructions. Their frames are not a match. None of their restrictions hold for bodily/abstract condition predications. We conclude by considering whether nominals other than those for body parts or abstractions can contribute to loss of the asymmetric force-dynamic relationship between participants that is central to the transitive prototype.
Key words: transitivity prominence, peripheral predications, Edoid, Emai}

\IfFileExists{../localcommands.tex}{
   \addbibresource{../localbibliography.bib}
   %%%%%%%%%%%%%%%%%%%%%%%%%%%%%%%%%%%%%%%%%%%%%
%%%%%%%%%% From LSP %%%%%%%%%%%%%%%%%%%%%%%%%
%%%%%%%%%%%%%%%%%%%%%%%%%%%%%%%%%%%%%%%%%%%%%

\usepackage{tabularx,multicol}
\usepackage{url}
\urlstyle{same}

\usepackage{langsci-optional}
\usepackage{langsci-lgr}
\usepackage{langsci-gb4e}

% NOTE THAT THE FOLLOWING PACKAGES ARE BANNED: 
% - pstricks
% - tipa
%%%%%%%%%%%%%%%%%%%%%%%%%%%%%%%%%%%%%%%%%%%%%
%%%%%%%%%% From ACAL LaTeX template %%%%%%%%%
%%%%%%%%%%%%%%%%%%%%%%%%%%%%%%%%%%%%%%%%%%%%%

%For stacking text, used here in autosegmental diagrams
\usepackage{stackengine}

%To combine rows in tables
\usepackage{multirow}

%Gives some extra formatting options, e.g. underlining/strikeout
\usepackage[normalem]{ulem}

%Use package/command below to create a double-spaced document, if you want one. Uncomment BOTH the package and the command (\doublespacing) to create a doublespaced document, or leave them as is to have a single-spaced document.
%\usepackage{setspace}
%\doublespacing 

 

%use for special OT tableaux symbols like bomb and sad face. must be loaded early on because it doesn't play well with some other packages
\usepackage{fourier-orns}

%Basic math symbols 
\usepackage{pifont}
\usepackage{amssymb}

%%%Gives shortcuts for glossing. The use of this package is NOT explained in the Quick Reference Guide, but the documentation is on CTAN for those that are interested. MJKD finds it handy for glossing. (https://ctan.org/pkg/leipzig?lang=en)
\usepackage{leipzig}

%Tables
% \usepackage{caption} %For table captions 

%For OT-style tableaux
\usepackage[notipa]{ot-tableau}

%highlights text with \hl{text}
\usepackage{color, soul} %note that soul sometimes eats Unicode text witouth telling you. 

%Drawing Syntax Trees
\usepackage[linguistics]{forest}

%This specifies some formatting for the forest trees to make them nicer to look at. Unfortunately this was borrowed by Michael Diercks from someone a while ago and he can't remember who. Apologies and if someone lets him know he will update this.
\forestset{
  nice nodes/.style={
    for tree={
      inner sep=0pt,
      fit=band,
    },
  },
  default preamble=nice nodes,
}

%A couple of packages that seemed to prefer being called toward the end of the preamble

%This package provides macros for typesetting SPE-style phonological rules. I've redefined the \oneof command here, so that there the rule includes a right brace. 
\usepackage{phonrule}
\renewcommand{\oneof}[2][c]{\ensuremath{
    ~\left\{
    \begin{tabular}{#1#1}#2\end{tabular}
    \right\}
    }}

%For using Greek letters outside of math mode.
% \usepackage{textgreek}


%Random, lets us use the XeLaTeX logo. Not important to the template at all.
% \usepackage{metalogo}



%%%%%%%%%%%%%%%%%%%%%%%%%%%%%%%%%%%%%%%%%%%%%
%%%%%%%%%% From Smith paper %%%%%%%%%%%%%%%%%
%%%%%%%%%%%%%%%%%%%%%%%%%%%%%%%%%%%%%%%%%%%%%

\usetikzlibrary{positioning}
\usetikzlibrary{trees}
% \usepackage{pstricks} %pstricks is banned. Use tikz instead. 
% \usepackage{pst-node}
%\usepackage{tikz}

%%%%%%%%%%%%%%%%%%%%%%%%%%%%%%%%%%%%%%%%%%%%%
%%%%%%%%%% From Gluckman paper %%%%%%%%%%%%%%
%%%%%%%%%%%%%%%%%%%%%%%%%%%%%%%%%%%%%%%%%%%%%

\usepackage{amsmath}



%%%%%%%%%%%%%%%%%%%%%%%%%%%%%%%%%%%%%%%%%%%%%
%%%%%%%%%% From Kandybowicz paper %%%%%%%%%%%
%%%%%%%%%%%%%%%%%%%%%%%%%%%%%%%%%%%%%%%%%%%%%

%These are in the .tex, but there's nothing in the "Figures" folder so I'm not sure that it's actually doing anything. 
 
% \graphicspath{ {./Figures/} }

%%%%%%%%%%%%%%%%%%%%%%%%%%%%%%%%%%%%%%%%%%%%%
%%%%%%%%%% From Matte paper %%%%%%%%%%%%%%%%%
%%%%%%%%%%%%%%%%%%%%%%%%%%%%%%%%%%%%%%%%%%%%%

%%%%%%%%%%%%%%%%%%%%%%%%%%%%%%%%%%%%%%%%%%%%%
%%%%%%%%%% From Grollemund paper %%%%%%%%%%%%%%
%%%%%%%%%%%%%%%%%%%%%%%%%%%%%%%%%%%%%%%%%%%%%

% \usepackage{lscape}

%%%%%%%%%%%%%%%%%%%%%%%%%%%%%%%%%%%%%%%%%%%%%
%%%%%%%%%% From Burns paper %%%%%%%%%%%%%%
%%%%%%%%%%%%%%%%%%%%%%%%%%%%%%%%%%%%%%%%%%%%%


% \usepackage{url}
% \urlstyle{same}

\usepackage{listings}
\lstset{basicstyle=\ttfamily,tabsize=2,breaklines=true}

%MJKD commented out - these are standard for chapters in edited volumes, I don't think we need them here in the preamble. 

% \usepackage{tabularx,multicol}
% \usepackage{langsci-basic}
% \usepackage{langsci-gb4e}
% \bibliography{localbibliography}
% \newcommand{\orcid}[1]{}
% \pagenumbering{roman}

% \IfFileExists{../localcommands.tex}{%hack to check whether this is being compiled as part of a collection or standalone
%    %%%%%%%%%%%%%%%%%%%%%%%%%%%%%%%%%%%%%%%%%%%%%
%%%%%%%%%% From LSP %%%%%%%%%%%%%%%%%%%%%%%%%
%%%%%%%%%%%%%%%%%%%%%%%%%%%%%%%%%%%%%%%%%%%%%

\usepackage{tabularx,multicol}
\usepackage{url}
\urlstyle{same}

\usepackage{langsci-optional}
\usepackage{langsci-lgr}
\usepackage{langsci-gb4e}

% NOTE THAT THE FOLLOWING PACKAGES ARE BANNED: 
% - pstricks
% - tipa
%%%%%%%%%%%%%%%%%%%%%%%%%%%%%%%%%%%%%%%%%%%%%
%%%%%%%%%% From ACAL LaTeX template %%%%%%%%%
%%%%%%%%%%%%%%%%%%%%%%%%%%%%%%%%%%%%%%%%%%%%%

%For stacking text, used here in autosegmental diagrams
\usepackage{stackengine}

%To combine rows in tables
\usepackage{multirow}

%Gives some extra formatting options, e.g. underlining/strikeout
\usepackage[normalem]{ulem}

%Use package/command below to create a double-spaced document, if you want one. Uncomment BOTH the package and the command (\doublespacing) to create a doublespaced document, or leave them as is to have a single-spaced document.
%\usepackage{setspace}
%\doublespacing 

 

%use for special OT tableaux symbols like bomb and sad face. must be loaded early on because it doesn't play well with some other packages
\usepackage{fourier-orns}

%Basic math symbols 
\usepackage{pifont}
\usepackage{amssymb}

%%%Gives shortcuts for glossing. The use of this package is NOT explained in the Quick Reference Guide, but the documentation is on CTAN for those that are interested. MJKD finds it handy for glossing. (https://ctan.org/pkg/leipzig?lang=en)
\usepackage{leipzig}

%Tables
% \usepackage{caption} %For table captions 

%For OT-style tableaux
\usepackage[notipa]{ot-tableau}

%highlights text with \hl{text}
\usepackage{color, soul} %note that soul sometimes eats Unicode text witouth telling you. 

%Drawing Syntax Trees
\usepackage[linguistics]{forest}

%This specifies some formatting for the forest trees to make them nicer to look at. Unfortunately this was borrowed by Michael Diercks from someone a while ago and he can't remember who. Apologies and if someone lets him know he will update this.
\forestset{
  nice nodes/.style={
    for tree={
      inner sep=0pt,
      fit=band,
    },
  },
  default preamble=nice nodes,
}

%A couple of packages that seemed to prefer being called toward the end of the preamble

%This package provides macros for typesetting SPE-style phonological rules. I've redefined the \oneof command here, so that there the rule includes a right brace. 
\usepackage{phonrule}
\renewcommand{\oneof}[2][c]{\ensuremath{
    ~\left\{
    \begin{tabular}{#1#1}#2\end{tabular}
    \right\}
    }}

%For using Greek letters outside of math mode.
% \usepackage{textgreek}


%Random, lets us use the XeLaTeX logo. Not important to the template at all.
% \usepackage{metalogo}



%%%%%%%%%%%%%%%%%%%%%%%%%%%%%%%%%%%%%%%%%%%%%
%%%%%%%%%% From Smith paper %%%%%%%%%%%%%%%%%
%%%%%%%%%%%%%%%%%%%%%%%%%%%%%%%%%%%%%%%%%%%%%

\usetikzlibrary{positioning}
\usetikzlibrary{trees}
% \usepackage{pstricks} %pstricks is banned. Use tikz instead. 
% \usepackage{pst-node}
%\usepackage{tikz}

%%%%%%%%%%%%%%%%%%%%%%%%%%%%%%%%%%%%%%%%%%%%%
%%%%%%%%%% From Gluckman paper %%%%%%%%%%%%%%
%%%%%%%%%%%%%%%%%%%%%%%%%%%%%%%%%%%%%%%%%%%%%

\usepackage{amsmath}



%%%%%%%%%%%%%%%%%%%%%%%%%%%%%%%%%%%%%%%%%%%%%
%%%%%%%%%% From Kandybowicz paper %%%%%%%%%%%
%%%%%%%%%%%%%%%%%%%%%%%%%%%%%%%%%%%%%%%%%%%%%

%These are in the .tex, but there's nothing in the "Figures" folder so I'm not sure that it's actually doing anything. 
 
% \graphicspath{ {./Figures/} }

%%%%%%%%%%%%%%%%%%%%%%%%%%%%%%%%%%%%%%%%%%%%%
%%%%%%%%%% From Matte paper %%%%%%%%%%%%%%%%%
%%%%%%%%%%%%%%%%%%%%%%%%%%%%%%%%%%%%%%%%%%%%%

%%%%%%%%%%%%%%%%%%%%%%%%%%%%%%%%%%%%%%%%%%%%%
%%%%%%%%%% From Grollemund paper %%%%%%%%%%%%%%
%%%%%%%%%%%%%%%%%%%%%%%%%%%%%%%%%%%%%%%%%%%%%

% \usepackage{lscape}

%%%%%%%%%%%%%%%%%%%%%%%%%%%%%%%%%%%%%%%%%%%%%
%%%%%%%%%% From Burns paper %%%%%%%%%%%%%%
%%%%%%%%%%%%%%%%%%%%%%%%%%%%%%%%%%%%%%%%%%%%%


% \usepackage{url}
% \urlstyle{same}

\usepackage{listings}
\lstset{basicstyle=\ttfamily,tabsize=2,breaklines=true}

%MJKD commented out - these are standard for chapters in edited volumes, I don't think we need them here in the preamble. 

% \usepackage{tabularx,multicol}
% \usepackage{langsci-basic}
% \usepackage{langsci-gb4e}
% \bibliography{localbibliography}
% \newcommand{\orcid}[1]{}
% \pagenumbering{roman}

% \IfFileExists{../localcommands.tex}{%hack to check whether this is being compiled as part of a collection or standalone
%    %%%%%%%%%%%%%%%%%%%%%%%%%%%%%%%%%%%%%%%%%%%%%
%%%%%%%%%% From LSP %%%%%%%%%%%%%%%%%%%%%%%%%
%%%%%%%%%%%%%%%%%%%%%%%%%%%%%%%%%%%%%%%%%%%%%

\usepackage{tabularx,multicol}
\usepackage{url}
\urlstyle{same}

\usepackage{langsci-optional}
\usepackage{langsci-lgr}
\usepackage{langsci-gb4e}

% NOTE THAT THE FOLLOWING PACKAGES ARE BANNED: 
% - pstricks
% - tipa
%%%%%%%%%%%%%%%%%%%%%%%%%%%%%%%%%%%%%%%%%%%%%
%%%%%%%%%% From ACAL LaTeX template %%%%%%%%%
%%%%%%%%%%%%%%%%%%%%%%%%%%%%%%%%%%%%%%%%%%%%%

%For stacking text, used here in autosegmental diagrams
\usepackage{stackengine}

%To combine rows in tables
\usepackage{multirow}

%Gives some extra formatting options, e.g. underlining/strikeout
\usepackage[normalem]{ulem}

%Use package/command below to create a double-spaced document, if you want one. Uncomment BOTH the package and the command (\doublespacing) to create a doublespaced document, or leave them as is to have a single-spaced document.
%\usepackage{setspace}
%\doublespacing 

 

%use for special OT tableaux symbols like bomb and sad face. must be loaded early on because it doesn't play well with some other packages
\usepackage{fourier-orns}

%Basic math symbols 
\usepackage{pifont}
\usepackage{amssymb}

%%%Gives shortcuts for glossing. The use of this package is NOT explained in the Quick Reference Guide, but the documentation is on CTAN for those that are interested. MJKD finds it handy for glossing. (https://ctan.org/pkg/leipzig?lang=en)
\usepackage{leipzig}

%Tables
% \usepackage{caption} %For table captions 

%For OT-style tableaux
\usepackage[notipa]{ot-tableau}

%highlights text with \hl{text}
\usepackage{color, soul} %note that soul sometimes eats Unicode text witouth telling you. 

%Drawing Syntax Trees
\usepackage[linguistics]{forest}

%This specifies some formatting for the forest trees to make them nicer to look at. Unfortunately this was borrowed by Michael Diercks from someone a while ago and he can't remember who. Apologies and if someone lets him know he will update this.
\forestset{
  nice nodes/.style={
    for tree={
      inner sep=0pt,
      fit=band,
    },
  },
  default preamble=nice nodes,
}

%A couple of packages that seemed to prefer being called toward the end of the preamble

%This package provides macros for typesetting SPE-style phonological rules. I've redefined the \oneof command here, so that there the rule includes a right brace. 
\usepackage{phonrule}
\renewcommand{\oneof}[2][c]{\ensuremath{
    ~\left\{
    \begin{tabular}{#1#1}#2\end{tabular}
    \right\}
    }}

%For using Greek letters outside of math mode.
% \usepackage{textgreek}


%Random, lets us use the XeLaTeX logo. Not important to the template at all.
% \usepackage{metalogo}



%%%%%%%%%%%%%%%%%%%%%%%%%%%%%%%%%%%%%%%%%%%%%
%%%%%%%%%% From Smith paper %%%%%%%%%%%%%%%%%
%%%%%%%%%%%%%%%%%%%%%%%%%%%%%%%%%%%%%%%%%%%%%

\usetikzlibrary{positioning}
\usetikzlibrary{trees}
% \usepackage{pstricks} %pstricks is banned. Use tikz instead. 
% \usepackage{pst-node}
%\usepackage{tikz}

%%%%%%%%%%%%%%%%%%%%%%%%%%%%%%%%%%%%%%%%%%%%%
%%%%%%%%%% From Gluckman paper %%%%%%%%%%%%%%
%%%%%%%%%%%%%%%%%%%%%%%%%%%%%%%%%%%%%%%%%%%%%

\usepackage{amsmath}



%%%%%%%%%%%%%%%%%%%%%%%%%%%%%%%%%%%%%%%%%%%%%
%%%%%%%%%% From Kandybowicz paper %%%%%%%%%%%
%%%%%%%%%%%%%%%%%%%%%%%%%%%%%%%%%%%%%%%%%%%%%

%These are in the .tex, but there's nothing in the "Figures" folder so I'm not sure that it's actually doing anything. 
 
% \graphicspath{ {./Figures/} }

%%%%%%%%%%%%%%%%%%%%%%%%%%%%%%%%%%%%%%%%%%%%%
%%%%%%%%%% From Matte paper %%%%%%%%%%%%%%%%%
%%%%%%%%%%%%%%%%%%%%%%%%%%%%%%%%%%%%%%%%%%%%%

%%%%%%%%%%%%%%%%%%%%%%%%%%%%%%%%%%%%%%%%%%%%%
%%%%%%%%%% From Grollemund paper %%%%%%%%%%%%%%
%%%%%%%%%%%%%%%%%%%%%%%%%%%%%%%%%%%%%%%%%%%%%

% \usepackage{lscape}

%%%%%%%%%%%%%%%%%%%%%%%%%%%%%%%%%%%%%%%%%%%%%
%%%%%%%%%% From Burns paper %%%%%%%%%%%%%%
%%%%%%%%%%%%%%%%%%%%%%%%%%%%%%%%%%%%%%%%%%%%%


% \usepackage{url}
% \urlstyle{same}

\usepackage{listings}
\lstset{basicstyle=\ttfamily,tabsize=2,breaklines=true}

%MJKD commented out - these are standard for chapters in edited volumes, I don't think we need them here in the preamble. 

% \usepackage{tabularx,multicol}
% \usepackage{langsci-basic}
% \usepackage{langsci-gb4e}
% \bibliography{localbibliography}
% \newcommand{\orcid}[1]{}
% \pagenumbering{roman}

% \IfFileExists{../localcommands.tex}{%hack to check whether this is being compiled as part of a collection or standalone
%    \input{../localpackages}
%    \input{../localcommands}
%    \input{../localhyphenation}
%     \bibliography{localbibliography}
%     \togglepaper[23]
% }{}

% %custom footer for preprints
% \papernote{\scriptsize\normalfont
%     Change author.
%     Change title.
%     To appear in:
%     Change Volume Editor.  
%     Change volume title.  
%     Berlin: Language Science Press. [preliminary page numbering]
% }



%%%%%%%%%%%%%%%%%%%%%%%%%%%%%%%%%%%%%%%%%%%%%
%%%%%%%%%% From Feibig paper %%%%%%%%%%%%%%
%%%%%%%%%%%%%%%%%%%%%%%%%%%%%%%%%%%%%%%%%%%%%

\usepackage{tikz-qtree}
\usepackage{tikz-qtree-compat}

%%%%%%%%%%%%%%%%%%%%%%%%%%%%%%%%%%%%%%%%%%%%%
%%%%%%%%%% From Olson paper %%%%%%%%%%%%%%
%%%%%%%%%%%%%%%%%%%%%%%%%%%%%%%%%%%%%%%%%%%%%

%MJKD - TIPA is forbidden lol. Commenting out for now but we'll have to address this in the Olson paper

%\usepackage[tone]{tipa}

%%%%%%%%%%%%%%%%%%%%%%%%%%%%%%%%%%%%%%%%%%%%%
%%%%%%%%%% From Caragine paper %%%%%%%%%%%%%%
%%%%%%%%%%%%%%%%%%%%%%%%%%%%%%%%%%%%%%%%%%%%%

\usepackage{placeins}
% \usepackage{graphicx}
% \usepackage{float} %Overleaf suggested this as a solution


%%%%%%%%%%%%%%%%%%%%%%%%%%%%%%%%%%%%%%%%%%%%%
%%%%%%%%%% From Kieffer paper %%%%%%%%%%%%%%
%%%%%%%%%%%%%%%%%%%%%%%%%%%%%%%%%%%%%%%%%%%%%

\usepackage{array}

%%%%%%%%%%%%%%%%%%%%%%%%%%%%%%%%%%%%%%%%%%%%%
%%%%%%%%%% From Payne paper %%%%%%%%%%%%%%
%%%%%%%%%%%%%%%%%%%%%%%%%%%%%%%%%%%%%%%%%%%%%

\usepackage{enumerate}
\usepackage{array,ragged2e}
\usepackage{pgfplots}


%%%%%%%%%%%%%%%%%%%%%%%%%%%%%%%%%%%%%%%%%%%%%
%%%%%%%%%% From Diercks paper %%%%%%%%%%%%%%
%%%%%%%%%%%%%%%%%%%%%%%%%%%%%%%%%%%%%%%%%%%%%

% \usepackage{cgloss}



%%%%%%%%%%%%%%%%%%%%%%%%%%%%%%%%%%%%%%%%%%%%%
%%%%%%%%%% From Li paper %%%%%%%%%%%%%%
%%%%%%%%%%%%%%%%%%%%%%%%%%%%%%%%%%%%%%%%%%%%%



%%%%%%%%%%%%%%%%%%%%%%%%%%%%%%%%%%%%%%%%%%%%%
%%%%%%%%%% From Myers paper %%%%%%%%%%%%%%
%%%%%%%%%%%%%%%%%%%%%%%%%%%%%%%%%%%%%%%%%%%%%



\usepackage{tikz-qtree}
\usepackage{tikz-qtree-compat}


%Package that makes glossing slightly easier.
\RequirePackage{leipzig}
%\RequirePackage{mathtools} % for \mathrlap

% % % Shortcuts (various borrowed from Jason Zentz)
\RequirePackage{xspace}
\xspaceaddexceptions{]\}}
\usepackage{langsci-textipa}
\usepackage{langsci-branding}
\usepackage{tabto}

%    %%%%%%%%%%%%%%%%%%%%%%%%%%%%%%%%%%%%%%%%%%%%%
%%%%%%%%%% From LSP %%%%%%%%%%%%%%%%%%%%%%%%%
%%%%%%%%%%%%%%%%%%%%%%%%%%%%%%%%%%%%%%%%%%%%%


% \newcommand*{\orcid}{}


\makeatletter
\let\thetitle\@title
\let\theauthor\@author
\makeatother

\newcommand{\togglepaper}[1][0]{
   \bibliography{../localbibliography}
   \papernote{\scriptsize\normalfont
     \theauthor.
     \titleTemp.
     To appear in:
     E. Di Tor \& Herr Rausgeberin (ed.).
     Booktitle in localcommands.tex.
     Berlin: Language Science Press. [preliminary page numbering]
   }
   \pagenumbering{roman}
   \setcounter{chapter}{#1}
   \addtocounter{chapter}{-1}
}

\newbool{bookcompile}
\booltrue{bookcompile}
\newcommand{\bookorchapter}[2]{\ifbool{bookcompile}{#1}{#2}}


%%%%%%%%%%%%%%%%%%%%%%%%%%%%%%%%%%%%%%%%%%%%%
%%%%%%%%%% From ACAL LaTeX template %%%%%%%%%
%%%%%%%%%%%%%%%%%%%%%%%%%%%%%%%%%%%%%%%%%%%%%


\usepackage{tikz}

\newcommand{\circled}[1]{\begin{tikzpicture}[baseline=(word.base)]
\node[draw, rounded corners, text height=8pt, text depth=2pt, inner sep=2pt, outer sep=0pt, use as bounding box] (word) {#1};
\end{tikzpicture}
}


%%%%%%%%%%%%%%%%%%%%%%%%%%%%%%%%%%%%%%%%%%%%%%
%%%%%%%%%%%%%%%%%%%%%%%%%%%%%%%%%%%%%%%%%%%%%%

% Useful Ling Shortcuts


%This makes the \emptyset command be a nicer one
\let\oldemptyset\emptyset
\let\emptyset\varnothing
\newcommand{\nothing}{$\emptyset$}

%Not all of these are explained in the Quick Reference Guide, but they are here bc they are relevant to some of our students. 
\newcommand{\xbar}{\ensuremath{'}\xspace}
\newcommand{\xhead}{\ensuremath{^\circ}\xspace}
\newcommand{\Lb}[1]{$\text{[}_{\text{#1}}$ } %A more convenient left bracket
\newcommand{\Rb}[1]{$\text{]}_{\text{#1}}$ } %A more convenient left bracket
\newcommand{\gap}{\underline{\hspace{1.2em}}}
\newcommand{\vP}{\emph{v}P}
\newcommand{\lilv}{\emph{v}}
\newcommand{\Abar}{A$'$-} %A more convenient A-bar notation
\newcommand{\ph}{$\varphi$\xspace} %A more convenient phi
\newcommand{\pro}{\emph{pro}\xspace}
\newcommand{\subs}[1]{\textsubscript{#1}} %A more convenient subscript
\newcommand{\spells}{$\Longleftrightarrow$} %spellout arrow for morph spellout rules
\newcommand{\tr}[1]{\textit{t}\textsubscript{\textit{#1}}} %easy traces with subscript
\newcommand{\supers}[1]{\textsuperscript{#1}}

%These are borrowed/modified from Andrew Mackenzie's ling-macros package. We have copied them in here (instead of just using the package itself) to avoid a minor clash that was generating an error.


% Abbreviations for glossing, based on Leipzig
\newleipzig{hab}{hab}{habitual}
\newleipzig{rem}{rem}{remote}
\newleipzig{sm}{sm}{subject marker}
\newleipzig{t}{t}{tense}
\newleipzig{aa}{aa}{anti-agreement}
\newleipzig{pron}{pron}{pronoun}
\newleipzig{rec}{rec}{recent}
\newleipzig{om}{om}{object marker}
%\newleipzig{ipfv}{ipfv}{imperfective}
\newleipzig{asp}{asp}{aspect}
\newleipzig{lk}{lk}{linker}
\newleipzig{pcl}{pcl}{particle}
\newleipzig{stat}{stat}{stative}
\newleipzig{ints}{ints}{intensive}
\newleipzig{ascl}{ascl}{assertive subject clitic}
\newleipzig{nascl}{nascl}{non-assertive subject clitic}
\newleipzig{ta}{ta}{tense and/or aspect}
\newleipzig{assoc}{assoc}{associative marker}
\newleipzig{hon}{hon}{honorific}
%\newleipzig{whprt}{wh}{\wh particle}
\newleipzig{sa}{sa}{subject agreement}
\newleipzig{conj}{conj}{conjunction}
%\newleipzig{loc}{loc}{locative}
\newleipzig{expl}{expl}{expletive}
\newleipzig{rcm}{rcm}{reciprocal marker}
\newleipzig{pers}{pers}{persistive}
\newleipzig{fv}{fv}{final vowel}
%\newleipzig{}{}{} %this is just to copy for when I want to add more


%%%%%%%%%%%%%%%%%%%%%%%%%%%%%%%%%%%%%%%%%%%%%%%%%%%%%
%%%%%%%%%%%%%%%%%%%%%%%%%%%%%%%%%%%%%%%%%%%%%%%%%%%%%


%%%%%%%%%%%%%%%%%%%%%%%%%%%%%%%%%%%%%%%%%%%%%
%%%%%%%%%% From Gluckman paper %%%%%%%%%%%%%%
%%%%%%%%%%%%%%%%%%%%%%%%%%%%%%%%%%%%%%%%%%%%%

\newcommand{\pcite}[1]{\citeauthor{#1}'s (\citeyear{#1})}


%%%%%%%%%%%%%%%%%%%%%%%%%%%%%%%%%%%%%%%%%%%%%
%%%%%%%%%% From Myers paper %%%%%%%%%%%%%%
%%%%%%%%%%%%%%%%%%%%%%%%%%%%%%%%%%%%%%%%%%%%%

%% hspace over width of something without showing it
\newcommand*{\hspaceThis}[1]{\settowidth{\LSPTmp}{#1}\hspace*{\LSPTmp}}
% use \hphantom{} for this


%%%%%%%%%%%%%%%%%%%%%%%%%%%%%%%%%%%%%%%%%%%%%
%%%%%%%%%% From Matte paper %%%%%%%%%%%%%%%%%
%%%%%%%%%%%%%%%%%%%%%%%%%%%%%%%%%%%%%%%%%%%%%

%mjkd commented this out in case it's breaking something right now. 
% \newcounter{xlistex}
% \newcommand{\footnoteex}{\stepcounter{xlistex}(\roman{xlistex})}

\newleipzig{sp}{sp}{speaker} %a new leipzig glossing command


%%%%%%%%%%%%%%%%%%%%%%%%%%%%%%%%%%%%%%%%%%%%%
%%%%%%%%%% From GamFranich paper %%%%%%%%%%%%%%
%%%%%%%%%%%%%%%%%%%%%%%%%%%%%%%%%%%%%%%%%%%%%

%MJKD commented these out - both will be provided by the overall LSP book template 

% \bibliography{localbibliography}
% \newcommand{\orcid}[1]{}

%%%%%%%%%%%%%%%%%%%%%%%%%%%%%%%%%%%%%%%%%%%%%
%%%%%%%%%% From Diercks paper %%%%%%%%%%%%%%
%%%%%%%%%%%%%%%%%%%%%%%%%%%%%%%%%%%%%%%%%%%%%

\newcommand{\abar}{$\bar{\text{A}}$\xspace} % this one is different to the \Abar command in Mike's shortcuts doc
\newcommand{\topFeat}{[\textsc{top}]\xspace}


%%%%%%%%%%%%%%%%%%%%%%%%%%%%%%%%%%%%%%%%%%%%%
%%%%%%%%%% From Kinchsular paper %%%%%%%%%%%%%%
%%%%%%%%%%%%%%%%%%%%%%%%%%%%%%%%%%%%%%%%%%%%%

%%%%%%%%%%%%%%%%%%%%%%%%%%%%%%%%%%%%%%%%%%%%%
%%%%%%%%%% From Chen paper %%%%%%%%%%%%%%
%%%%%%%%%%%%%%%%%%%%%%%%%%%%%%%%%%%%%%%%%%%%%

\newcounter{savedex}
\makeatletter
\newcommand*{\saveex}{\setcounter{savedex}{\value{\@xnumctr}}}
\newcommand*{\resumeex}{\setcounter{\@xnumctr}{\value{savedex}}}
\makeatother


%    \input{../localhyphenation}
%     \bibliography{localbibliography}
%     \togglepaper[23]
% }{}

% %custom footer for preprints
% \papernote{\scriptsize\normalfont
%     Change author.
%     Change title.
%     To appear in:
%     Change Volume Editor.  
%     Change volume title.  
%     Berlin: Language Science Press. [preliminary page numbering]
% }



%%%%%%%%%%%%%%%%%%%%%%%%%%%%%%%%%%%%%%%%%%%%%
%%%%%%%%%% From Feibig paper %%%%%%%%%%%%%%
%%%%%%%%%%%%%%%%%%%%%%%%%%%%%%%%%%%%%%%%%%%%%

\usepackage{tikz-qtree}
\usepackage{tikz-qtree-compat}

%%%%%%%%%%%%%%%%%%%%%%%%%%%%%%%%%%%%%%%%%%%%%
%%%%%%%%%% From Olson paper %%%%%%%%%%%%%%
%%%%%%%%%%%%%%%%%%%%%%%%%%%%%%%%%%%%%%%%%%%%%

%MJKD - TIPA is forbidden lol. Commenting out for now but we'll have to address this in the Olson paper

%\usepackage[tone]{tipa}

%%%%%%%%%%%%%%%%%%%%%%%%%%%%%%%%%%%%%%%%%%%%%
%%%%%%%%%% From Caragine paper %%%%%%%%%%%%%%
%%%%%%%%%%%%%%%%%%%%%%%%%%%%%%%%%%%%%%%%%%%%%

\usepackage{placeins}
% \usepackage{graphicx}
% \usepackage{float} %Overleaf suggested this as a solution


%%%%%%%%%%%%%%%%%%%%%%%%%%%%%%%%%%%%%%%%%%%%%
%%%%%%%%%% From Kieffer paper %%%%%%%%%%%%%%
%%%%%%%%%%%%%%%%%%%%%%%%%%%%%%%%%%%%%%%%%%%%%

\usepackage{array}

%%%%%%%%%%%%%%%%%%%%%%%%%%%%%%%%%%%%%%%%%%%%%
%%%%%%%%%% From Payne paper %%%%%%%%%%%%%%
%%%%%%%%%%%%%%%%%%%%%%%%%%%%%%%%%%%%%%%%%%%%%

\usepackage{enumerate}
\usepackage{array,ragged2e}
\usepackage{pgfplots}


%%%%%%%%%%%%%%%%%%%%%%%%%%%%%%%%%%%%%%%%%%%%%
%%%%%%%%%% From Diercks paper %%%%%%%%%%%%%%
%%%%%%%%%%%%%%%%%%%%%%%%%%%%%%%%%%%%%%%%%%%%%

% \usepackage{cgloss}



%%%%%%%%%%%%%%%%%%%%%%%%%%%%%%%%%%%%%%%%%%%%%
%%%%%%%%%% From Li paper %%%%%%%%%%%%%%
%%%%%%%%%%%%%%%%%%%%%%%%%%%%%%%%%%%%%%%%%%%%%



%%%%%%%%%%%%%%%%%%%%%%%%%%%%%%%%%%%%%%%%%%%%%
%%%%%%%%%% From Myers paper %%%%%%%%%%%%%%
%%%%%%%%%%%%%%%%%%%%%%%%%%%%%%%%%%%%%%%%%%%%%



\usepackage{tikz-qtree}
\usepackage{tikz-qtree-compat}


%Package that makes glossing slightly easier.
\RequirePackage{leipzig}
%\RequirePackage{mathtools} % for \mathrlap

% % % Shortcuts (various borrowed from Jason Zentz)
\RequirePackage{xspace}
\xspaceaddexceptions{]\}}
\usepackage{langsci-textipa}
\usepackage{langsci-branding}
\usepackage{tabto}

%    %%%%%%%%%%%%%%%%%%%%%%%%%%%%%%%%%%%%%%%%%%%%%
%%%%%%%%%% From LSP %%%%%%%%%%%%%%%%%%%%%%%%%
%%%%%%%%%%%%%%%%%%%%%%%%%%%%%%%%%%%%%%%%%%%%%


% \newcommand*{\orcid}{}


\makeatletter
\let\thetitle\@title
\let\theauthor\@author
\makeatother

\newcommand{\togglepaper}[1][0]{
   \bibliography{../localbibliography}
   \papernote{\scriptsize\normalfont
     \theauthor.
     \titleTemp.
     To appear in:
     E. Di Tor \& Herr Rausgeberin (ed.).
     Booktitle in localcommands.tex.
     Berlin: Language Science Press. [preliminary page numbering]
   }
   \pagenumbering{roman}
   \setcounter{chapter}{#1}
   \addtocounter{chapter}{-1}
}

\newbool{bookcompile}
\booltrue{bookcompile}
\newcommand{\bookorchapter}[2]{\ifbool{bookcompile}{#1}{#2}}


%%%%%%%%%%%%%%%%%%%%%%%%%%%%%%%%%%%%%%%%%%%%%
%%%%%%%%%% From ACAL LaTeX template %%%%%%%%%
%%%%%%%%%%%%%%%%%%%%%%%%%%%%%%%%%%%%%%%%%%%%%


\usepackage{tikz}

\newcommand{\circled}[1]{\begin{tikzpicture}[baseline=(word.base)]
\node[draw, rounded corners, text height=8pt, text depth=2pt, inner sep=2pt, outer sep=0pt, use as bounding box] (word) {#1};
\end{tikzpicture}
}


%%%%%%%%%%%%%%%%%%%%%%%%%%%%%%%%%%%%%%%%%%%%%%
%%%%%%%%%%%%%%%%%%%%%%%%%%%%%%%%%%%%%%%%%%%%%%

% Useful Ling Shortcuts


%This makes the \emptyset command be a nicer one
\let\oldemptyset\emptyset
\let\emptyset\varnothing
\newcommand{\nothing}{$\emptyset$}

%Not all of these are explained in the Quick Reference Guide, but they are here bc they are relevant to some of our students. 
\newcommand{\xbar}{\ensuremath{'}\xspace}
\newcommand{\xhead}{\ensuremath{^\circ}\xspace}
\newcommand{\Lb}[1]{$\text{[}_{\text{#1}}$ } %A more convenient left bracket
\newcommand{\Rb}[1]{$\text{]}_{\text{#1}}$ } %A more convenient left bracket
\newcommand{\gap}{\underline{\hspace{1.2em}}}
\newcommand{\vP}{\emph{v}P}
\newcommand{\lilv}{\emph{v}}
\newcommand{\Abar}{A$'$-} %A more convenient A-bar notation
\newcommand{\ph}{$\varphi$\xspace} %A more convenient phi
\newcommand{\pro}{\emph{pro}\xspace}
\newcommand{\subs}[1]{\textsubscript{#1}} %A more convenient subscript
\newcommand{\spells}{$\Longleftrightarrow$} %spellout arrow for morph spellout rules
\newcommand{\tr}[1]{\textit{t}\textsubscript{\textit{#1}}} %easy traces with subscript
\newcommand{\supers}[1]{\textsuperscript{#1}}

%These are borrowed/modified from Andrew Mackenzie's ling-macros package. We have copied them in here (instead of just using the package itself) to avoid a minor clash that was generating an error.


% Abbreviations for glossing, based on Leipzig
\newleipzig{hab}{hab}{habitual}
\newleipzig{rem}{rem}{remote}
\newleipzig{sm}{sm}{subject marker}
\newleipzig{t}{t}{tense}
\newleipzig{aa}{aa}{anti-agreement}
\newleipzig{pron}{pron}{pronoun}
\newleipzig{rec}{rec}{recent}
\newleipzig{om}{om}{object marker}
%\newleipzig{ipfv}{ipfv}{imperfective}
\newleipzig{asp}{asp}{aspect}
\newleipzig{lk}{lk}{linker}
\newleipzig{pcl}{pcl}{particle}
\newleipzig{stat}{stat}{stative}
\newleipzig{ints}{ints}{intensive}
\newleipzig{ascl}{ascl}{assertive subject clitic}
\newleipzig{nascl}{nascl}{non-assertive subject clitic}
\newleipzig{ta}{ta}{tense and/or aspect}
\newleipzig{assoc}{assoc}{associative marker}
\newleipzig{hon}{hon}{honorific}
%\newleipzig{whprt}{wh}{\wh particle}
\newleipzig{sa}{sa}{subject agreement}
\newleipzig{conj}{conj}{conjunction}
%\newleipzig{loc}{loc}{locative}
\newleipzig{expl}{expl}{expletive}
\newleipzig{rcm}{rcm}{reciprocal marker}
\newleipzig{pers}{pers}{persistive}
\newleipzig{fv}{fv}{final vowel}
%\newleipzig{}{}{} %this is just to copy for when I want to add more


%%%%%%%%%%%%%%%%%%%%%%%%%%%%%%%%%%%%%%%%%%%%%%%%%%%%%
%%%%%%%%%%%%%%%%%%%%%%%%%%%%%%%%%%%%%%%%%%%%%%%%%%%%%


%%%%%%%%%%%%%%%%%%%%%%%%%%%%%%%%%%%%%%%%%%%%%
%%%%%%%%%% From Gluckman paper %%%%%%%%%%%%%%
%%%%%%%%%%%%%%%%%%%%%%%%%%%%%%%%%%%%%%%%%%%%%

\newcommand{\pcite}[1]{\citeauthor{#1}'s (\citeyear{#1})}


%%%%%%%%%%%%%%%%%%%%%%%%%%%%%%%%%%%%%%%%%%%%%
%%%%%%%%%% From Myers paper %%%%%%%%%%%%%%
%%%%%%%%%%%%%%%%%%%%%%%%%%%%%%%%%%%%%%%%%%%%%

%% hspace over width of something without showing it
\newcommand*{\hspaceThis}[1]{\settowidth{\LSPTmp}{#1}\hspace*{\LSPTmp}}
% use \hphantom{} for this


%%%%%%%%%%%%%%%%%%%%%%%%%%%%%%%%%%%%%%%%%%%%%
%%%%%%%%%% From Matte paper %%%%%%%%%%%%%%%%%
%%%%%%%%%%%%%%%%%%%%%%%%%%%%%%%%%%%%%%%%%%%%%

%mjkd commented this out in case it's breaking something right now. 
% \newcounter{xlistex}
% \newcommand{\footnoteex}{\stepcounter{xlistex}(\roman{xlistex})}

\newleipzig{sp}{sp}{speaker} %a new leipzig glossing command


%%%%%%%%%%%%%%%%%%%%%%%%%%%%%%%%%%%%%%%%%%%%%
%%%%%%%%%% From GamFranich paper %%%%%%%%%%%%%%
%%%%%%%%%%%%%%%%%%%%%%%%%%%%%%%%%%%%%%%%%%%%%

%MJKD commented these out - both will be provided by the overall LSP book template 

% \bibliography{localbibliography}
% \newcommand{\orcid}[1]{}

%%%%%%%%%%%%%%%%%%%%%%%%%%%%%%%%%%%%%%%%%%%%%
%%%%%%%%%% From Diercks paper %%%%%%%%%%%%%%
%%%%%%%%%%%%%%%%%%%%%%%%%%%%%%%%%%%%%%%%%%%%%

\newcommand{\abar}{$\bar{\text{A}}$\xspace} % this one is different to the \Abar command in Mike's shortcuts doc
\newcommand{\topFeat}{[\textsc{top}]\xspace}


%%%%%%%%%%%%%%%%%%%%%%%%%%%%%%%%%%%%%%%%%%%%%
%%%%%%%%%% From Kinchsular paper %%%%%%%%%%%%%%
%%%%%%%%%%%%%%%%%%%%%%%%%%%%%%%%%%%%%%%%%%%%%

%%%%%%%%%%%%%%%%%%%%%%%%%%%%%%%%%%%%%%%%%%%%%
%%%%%%%%%% From Chen paper %%%%%%%%%%%%%%
%%%%%%%%%%%%%%%%%%%%%%%%%%%%%%%%%%%%%%%%%%%%%

\newcounter{savedex}
\makeatletter
\newcommand*{\saveex}{\setcounter{savedex}{\value{\@xnumctr}}}
\newcommand*{\resumeex}{\setcounter{\@xnumctr}{\value{savedex}}}
\makeatother


%    \input{../localhyphenation}
%     \bibliography{localbibliography}
%     \togglepaper[23]
% }{}

% %custom footer for preprints
% \papernote{\scriptsize\normalfont
%     Change author.
%     Change title.
%     To appear in:
%     Change Volume Editor.  
%     Change volume title.  
%     Berlin: Language Science Press. [preliminary page numbering]
% }



%%%%%%%%%%%%%%%%%%%%%%%%%%%%%%%%%%%%%%%%%%%%%
%%%%%%%%%% From Feibig paper %%%%%%%%%%%%%%
%%%%%%%%%%%%%%%%%%%%%%%%%%%%%%%%%%%%%%%%%%%%%

\usepackage{tikz-qtree}
\usepackage{tikz-qtree-compat}

%%%%%%%%%%%%%%%%%%%%%%%%%%%%%%%%%%%%%%%%%%%%%
%%%%%%%%%% From Olson paper %%%%%%%%%%%%%%
%%%%%%%%%%%%%%%%%%%%%%%%%%%%%%%%%%%%%%%%%%%%%

%MJKD - TIPA is forbidden lol. Commenting out for now but we'll have to address this in the Olson paper

%\usepackage[tone]{tipa}

%%%%%%%%%%%%%%%%%%%%%%%%%%%%%%%%%%%%%%%%%%%%%
%%%%%%%%%% From Caragine paper %%%%%%%%%%%%%%
%%%%%%%%%%%%%%%%%%%%%%%%%%%%%%%%%%%%%%%%%%%%%

\usepackage{placeins}
% \usepackage{graphicx}
% \usepackage{float} %Overleaf suggested this as a solution


%%%%%%%%%%%%%%%%%%%%%%%%%%%%%%%%%%%%%%%%%%%%%
%%%%%%%%%% From Kieffer paper %%%%%%%%%%%%%%
%%%%%%%%%%%%%%%%%%%%%%%%%%%%%%%%%%%%%%%%%%%%%

\usepackage{array}

%%%%%%%%%%%%%%%%%%%%%%%%%%%%%%%%%%%%%%%%%%%%%
%%%%%%%%%% From Payne paper %%%%%%%%%%%%%%
%%%%%%%%%%%%%%%%%%%%%%%%%%%%%%%%%%%%%%%%%%%%%

\usepackage{enumerate}
\usepackage{array,ragged2e}
\usepackage{pgfplots}


%%%%%%%%%%%%%%%%%%%%%%%%%%%%%%%%%%%%%%%%%%%%%
%%%%%%%%%% From Diercks paper %%%%%%%%%%%%%%
%%%%%%%%%%%%%%%%%%%%%%%%%%%%%%%%%%%%%%%%%%%%%

% \usepackage{cgloss}



%%%%%%%%%%%%%%%%%%%%%%%%%%%%%%%%%%%%%%%%%%%%%
%%%%%%%%%% From Li paper %%%%%%%%%%%%%%
%%%%%%%%%%%%%%%%%%%%%%%%%%%%%%%%%%%%%%%%%%%%%



%%%%%%%%%%%%%%%%%%%%%%%%%%%%%%%%%%%%%%%%%%%%%
%%%%%%%%%% From Myers paper %%%%%%%%%%%%%%
%%%%%%%%%%%%%%%%%%%%%%%%%%%%%%%%%%%%%%%%%%%%%



\usepackage{tikz-qtree}
\usepackage{tikz-qtree-compat}


%Package that makes glossing slightly easier.
\RequirePackage{leipzig}
%\RequirePackage{mathtools} % for \mathrlap

% % % Shortcuts (various borrowed from Jason Zentz)
\RequirePackage{xspace}
\xspaceaddexceptions{]\}}
\usepackage{langsci-textipa}
\usepackage{langsci-branding}
\usepackage{tabto}

   %%%%%%%%%%%%%%%%%%%%%%%%%%%%%%%%%%%%%%%%%%%%%
%%%%%%%%%% From LSP %%%%%%%%%%%%%%%%%%%%%%%%%
%%%%%%%%%%%%%%%%%%%%%%%%%%%%%%%%%%%%%%%%%%%%%


% \newcommand*{\orcid}{}


\makeatletter
\let\thetitle\@title
\let\theauthor\@author
\makeatother

\newcommand{\togglepaper}[1][0]{
   \bibliography{../localbibliography}
   \papernote{\scriptsize\normalfont
     \theauthor.
     \titleTemp.
     To appear in:
     E. Di Tor \& Herr Rausgeberin (ed.).
     Booktitle in localcommands.tex.
     Berlin: Language Science Press. [preliminary page numbering]
   }
   \pagenumbering{roman}
   \setcounter{chapter}{#1}
   \addtocounter{chapter}{-1}
}

\newbool{bookcompile}
\booltrue{bookcompile}
\newcommand{\bookorchapter}[2]{\ifbool{bookcompile}{#1}{#2}}


%%%%%%%%%%%%%%%%%%%%%%%%%%%%%%%%%%%%%%%%%%%%%
%%%%%%%%%% From ACAL LaTeX template %%%%%%%%%
%%%%%%%%%%%%%%%%%%%%%%%%%%%%%%%%%%%%%%%%%%%%%


\usepackage{tikz}

\newcommand{\circled}[1]{\begin{tikzpicture}[baseline=(word.base)]
\node[draw, rounded corners, text height=8pt, text depth=2pt, inner sep=2pt, outer sep=0pt, use as bounding box] (word) {#1};
\end{tikzpicture}
}


%%%%%%%%%%%%%%%%%%%%%%%%%%%%%%%%%%%%%%%%%%%%%%
%%%%%%%%%%%%%%%%%%%%%%%%%%%%%%%%%%%%%%%%%%%%%%

% Useful Ling Shortcuts


%This makes the \emptyset command be a nicer one
\let\oldemptyset\emptyset
\let\emptyset\varnothing
\newcommand{\nothing}{$\emptyset$}

%Not all of these are explained in the Quick Reference Guide, but they are here bc they are relevant to some of our students. 
\newcommand{\xbar}{\ensuremath{'}\xspace}
\newcommand{\xhead}{\ensuremath{^\circ}\xspace}
\newcommand{\Lb}[1]{$\text{[}_{\text{#1}}$ } %A more convenient left bracket
\newcommand{\Rb}[1]{$\text{]}_{\text{#1}}$ } %A more convenient left bracket
\newcommand{\gap}{\underline{\hspace{1.2em}}}
\newcommand{\vP}{\emph{v}P}
\newcommand{\lilv}{\emph{v}}
\newcommand{\Abar}{A$'$-} %A more convenient A-bar notation
\newcommand{\ph}{$\varphi$\xspace} %A more convenient phi
\newcommand{\pro}{\emph{pro}\xspace}
\newcommand{\subs}[1]{\textsubscript{#1}} %A more convenient subscript
\newcommand{\spells}{$\Longleftrightarrow$} %spellout arrow for morph spellout rules
\newcommand{\tr}[1]{\textit{t}\textsubscript{\textit{#1}}} %easy traces with subscript
\newcommand{\supers}[1]{\textsuperscript{#1}}

%These are borrowed/modified from Andrew Mackenzie's ling-macros package. We have copied them in here (instead of just using the package itself) to avoid a minor clash that was generating an error.


% Abbreviations for glossing, based on Leipzig
\newleipzig{hab}{hab}{habitual}
\newleipzig{rem}{rem}{remote}
\newleipzig{sm}{sm}{subject marker}
\newleipzig{t}{t}{tense}
\newleipzig{aa}{aa}{anti-agreement}
\newleipzig{pron}{pron}{pronoun}
\newleipzig{rec}{rec}{recent}
\newleipzig{om}{om}{object marker}
%\newleipzig{ipfv}{ipfv}{imperfective}
\newleipzig{asp}{asp}{aspect}
\newleipzig{lk}{lk}{linker}
\newleipzig{pcl}{pcl}{particle}
\newleipzig{stat}{stat}{stative}
\newleipzig{ints}{ints}{intensive}
\newleipzig{ascl}{ascl}{assertive subject clitic}
\newleipzig{nascl}{nascl}{non-assertive subject clitic}
\newleipzig{ta}{ta}{tense and/or aspect}
\newleipzig{assoc}{assoc}{associative marker}
\newleipzig{hon}{hon}{honorific}
%\newleipzig{whprt}{wh}{\wh particle}
\newleipzig{sa}{sa}{subject agreement}
\newleipzig{conj}{conj}{conjunction}
%\newleipzig{loc}{loc}{locative}
\newleipzig{expl}{expl}{expletive}
\newleipzig{rcm}{rcm}{reciprocal marker}
\newleipzig{pers}{pers}{persistive}
\newleipzig{fv}{fv}{final vowel}
%\newleipzig{}{}{} %this is just to copy for when I want to add more


%%%%%%%%%%%%%%%%%%%%%%%%%%%%%%%%%%%%%%%%%%%%%%%%%%%%%
%%%%%%%%%%%%%%%%%%%%%%%%%%%%%%%%%%%%%%%%%%%%%%%%%%%%%


%%%%%%%%%%%%%%%%%%%%%%%%%%%%%%%%%%%%%%%%%%%%%
%%%%%%%%%% From Gluckman paper %%%%%%%%%%%%%%
%%%%%%%%%%%%%%%%%%%%%%%%%%%%%%%%%%%%%%%%%%%%%

\newcommand{\pcite}[1]{\citeauthor{#1}'s (\citeyear{#1})}


%%%%%%%%%%%%%%%%%%%%%%%%%%%%%%%%%%%%%%%%%%%%%
%%%%%%%%%% From Myers paper %%%%%%%%%%%%%%
%%%%%%%%%%%%%%%%%%%%%%%%%%%%%%%%%%%%%%%%%%%%%

%% hspace over width of something without showing it
\newcommand*{\hspaceThis}[1]{\settowidth{\LSPTmp}{#1}\hspace*{\LSPTmp}}
% use \hphantom{} for this


%%%%%%%%%%%%%%%%%%%%%%%%%%%%%%%%%%%%%%%%%%%%%
%%%%%%%%%% From Matte paper %%%%%%%%%%%%%%%%%
%%%%%%%%%%%%%%%%%%%%%%%%%%%%%%%%%%%%%%%%%%%%%

%mjkd commented this out in case it's breaking something right now. 
% \newcounter{xlistex}
% \newcommand{\footnoteex}{\stepcounter{xlistex}(\roman{xlistex})}

\newleipzig{sp}{sp}{speaker} %a new leipzig glossing command


%%%%%%%%%%%%%%%%%%%%%%%%%%%%%%%%%%%%%%%%%%%%%
%%%%%%%%%% From GamFranich paper %%%%%%%%%%%%%%
%%%%%%%%%%%%%%%%%%%%%%%%%%%%%%%%%%%%%%%%%%%%%

%MJKD commented these out - both will be provided by the overall LSP book template 

% \bibliography{localbibliography}
% \newcommand{\orcid}[1]{}

%%%%%%%%%%%%%%%%%%%%%%%%%%%%%%%%%%%%%%%%%%%%%
%%%%%%%%%% From Diercks paper %%%%%%%%%%%%%%
%%%%%%%%%%%%%%%%%%%%%%%%%%%%%%%%%%%%%%%%%%%%%

\newcommand{\abar}{$\bar{\text{A}}$\xspace} % this one is different to the \Abar command in Mike's shortcuts doc
\newcommand{\topFeat}{[\textsc{top}]\xspace}


%%%%%%%%%%%%%%%%%%%%%%%%%%%%%%%%%%%%%%%%%%%%%
%%%%%%%%%% From Kinchsular paper %%%%%%%%%%%%%%
%%%%%%%%%%%%%%%%%%%%%%%%%%%%%%%%%%%%%%%%%%%%%

%%%%%%%%%%%%%%%%%%%%%%%%%%%%%%%%%%%%%%%%%%%%%
%%%%%%%%%% From Chen paper %%%%%%%%%%%%%%
%%%%%%%%%%%%%%%%%%%%%%%%%%%%%%%%%%%%%%%%%%%%%

\newcounter{savedex}
\makeatletter
\newcommand*{\saveex}{\setcounter{savedex}{\value{\@xnumctr}}}
\newcommand*{\resumeex}{\setcounter{\@xnumctr}{\value{savedex}}}
\makeatother


   \input{../localhyphenation}
   \boolfalse{bookcompile}
   \togglepaper[23]%%chapternumber
}{}

\begin{document}
\maketitle

\section{Introduction}

The linguistic concept of transitivity remains underexplored despite a wealth of studies. At its prototypical core, transitivity has received rather extensive attention \citep{HopperandThompson1980, Tsunoda1985, Rice1987, Kittila2002, Kittila2008, Naess2007, Haspelmath2015}. The periphery of the transitivity concept, however, has been examined much less. This is especially so for languages in West Africa, where “hyper-transitivity” has been asserted \citep{AmekaEssegbey2006, Ameka2013, MuyskinSmith2015} but not defined or located relative to verb forms and their functions nor to nominal forms and their contribution to transitivity.  

For this paper we examine features of transitivity prominence in Emai, not at its core but at its periphery. Emai, an Edoid language of Nigeria \citet{Elugbe1989}, is SVO with lexical and grammatical tone, the latter relying on floating tones that are either high or low. It displays verbs that combine in series or that permit either postverbal aspectual particles or postverbal prepositions. Although these complex predicate types occur, we direct attention to simplex predicates, i.e. simple verb constructions, that assume a transitive shape, employ two arguments, and express a condition pertaining to humans and their bodies.

\section{Background}
\citet{Haspelmath2015} suggests that the perception of transitivity prominence may be more than impressionistic. He conducted a quantitative assessment of transitivity use, following a path initially explored with small verb sets and language sets by \citet{MullerGotama1994} and \citet{Bossong1998}. 

To assess transitivity prominence, Haspelmath sought to satisfy three criteria: 

\begin{itemize}
        \item A systematic database that would be reflective of languages from around the world;
        \item Verbs that would be diverse in meanings expressed and arguments used yet reasonably frequent in language use and therefore expected in most societies; 
        \item A rigorous definition of transitivity. 
\end{itemize}

These criteria were met in the Valency Patterns Leipzig electronic-database of \citet{HartmannHaspelmathTaylor2013}. The ValPaL database had a relatively large sample of verbs from thirty-five languages found around the globe. Verbs were defined as transitive if, like English ‘break,’ they had Agent and Patient arguments that correlated with the micro-roles “breaker” and “broken thing.” Eighty verb meanings were documented in ValPaL; their meanings were diverse in kind, conveyed event experiences common to human societies, and appeared to be reasonably frequent in everyday usage.   

Haspelmath employed a single measure of transitivity prominence across the ValPaL database. He computed the percentage of verbs in ValPaL that showed a transitive encoding of their meanings. Contrary to conventional wisdom about transitivity prominence in Central European languages, Haspelmath found that there were languages in Asia and Africa that had higher or much higher transitivity prominence. Some representative percentages from Haspelmath include the following: 

 \begin{itemize}
        \item Europe: Bezhta= .46; Russian= .56; English and German= .60; Italian= .64;
        \item Asia: Mandarin= .50; Jaminjung= .62; Japanese= .65; Balinese= .68; Chintang= .77;  
        \item Africa: Yoruba= .54; Standard Arabic= .62; N‖uu= .64; Mandinka= .64; Emai= .73. 
\end{itemize}

Among European languages, English and German had percentage scores somewhat above 50\%, whereas in Asia and Africa scores were found that exceeded 70\%, which no European language exhibited. Included among the percentage scores from Africa were those of Emai (73\%). Other languages from Africa did not have such a high score for the same set of meanings, e.g. Yoruba at 54\%. At least for the verbs and languages compared by Haspelmath, it appeared that a significant number of meanings in Emai employed a transitive coding and that such a number was among the highest recorded. In large part we undertake the following study to explore further the nature and scope of transitivity prominence in Emai.  

As for the set of meanings in English that stimulated the ValPaL database from thirty-five languages, they are presented in \tabref{tab:schaefer:frequencies}, taken from Haspelmath. Even in English, many of these meanings are coded by transitive verbs, but a number are coded with an intransitive form or a copula.  

\begin{table}
\caption{The 80 verb meanings of ValPaL}
\label{tab:schaefer:frequencies}
 \begin{tabularx}{.6\textwidth}{llll}
  \lsptoprule
eat  &   tell  &    put  &    die\\
hug  &   say  &    pour  &    play\\
look at  &   name  &    cover  &    be sad\\
see  &   build   &    fill  &    be hungry\\
smell  &  break   &    load  &    roll (intr.)\\
fear  &  kill    &    blink  &    sink (intr.)\\
frighten  &  beat   &    cough  &    burn (intr.)\\
like  &   hit  &    climb  &    be dry\\
know  &   touch  &    run  &    rain\\
think  &   cut  &    sit  &    be a hunter\\
search for  &   take  &    sit down  &    grind\\
wash  &   tear  &    jump  &    wipe\\
dress (tr.)  &   peel  &    sing  &    dig\\
shave  &   hide  &    go  &    push\\
help  &   show  &    leave  &    bring\\
follow  &   give  &    live  &    steal\\
meet  &   send  &    laugh  &    teach\\
talk  &   carry  &    scream  &    hear\\
ask for  &   throw  &    feel pain  &    cook\\
shout at  &   tie  &    feel cold  &    boil (intr.)\\
  \lspbottomrule
 \end{tabularx}
\end{table}

\section{Periphery of transitivity prominence in Emai}

There is more to transitivity and its prominence than the rather nebulous characterization “hyper-transitivity” that is found in the Africanist literature. Transitivity can be examined from the perspective of not only its core or prototype but also its periphery. But how does one explore the character of transitivity at its periphery? Where does one begin? Is there a domain that might be particularly revealing about the peripheral character of transitivity?  

Over the last decade or so a series of publications have investigated the human body and its coding by linguistic means (e.g. \cite{Ruthrof2015}, \cite{ZariquleyValenzuela2022}). Among them is \citet{BrenzingerKraskaSzlenk2014} with its focus on Africa. These studies address linguistic coding of the human body primarily by nominal forms; verbal forms that employ a body-part participant and encode a condition pertaining to a human participant have received far less attention.  

In the following we direct attention to two-participant predications that express experiential conditions applicable to humans and their bodily makeup. Subsequently, we consider two-participant predications that express activities of motion and bodily emissions. Both derive from fieldwork of various kinds with Emai, an Edoid language of central Nigeria. Through translation of Emai into English in various projects, we have compiled an extensive comparative corpus, including spontaneous texts from oral tradition and databases from directed elicitation fostering dictionary and grammar activities. 

Relative to illustrations from English, Emai bodily-condition predications deviate from the transitive prototype, as articulated by \citet{Naess2007}. Her principle of Maximally Distinguished Arguments (\textsc{mda}) holds that a canonical transitive clause exhibits two participants whose meanings are maximally distinct. Essentially, NP-1 is a volitional instigator, whereas NP-2 is a non-volitional, affected entity that registers a material effect. 

Deviation from the transitive prototype in Emai is evident in three primary predication types. At one level of analysis, these types differ with respect to the nature of their participants, which can be a human noun, an abstract noun, or a body-part noun. At a second level, they differ with respect to the positioning of the participant types. NP-1 can be a human noun or an abstract noun. NP-2 can be either a human noun, an abstract noun, or a body-part noun. Regardless of nominal nature or order, verb constructions expressing bodily-conditions constitute a two-participant predication (\textsc{tpp}), where a human participant occupies one position and a body-part noun or abstract noun the other position. Deviation from the prototype is evident in the failure of participant distinctness and the failure of material effect rather than simple overuse of a two-place predicate. 

 \textsc{tpp}s of one subtype diverge from the \textsc{mda} principle based on the participant type in NP-2 position. The examples from Emai in (\ref{ex:schaefer:sheg1}) show that NP-1 is a human noun and NP-2 is an abstract noun that conveys a physical or non-physical bodily condition. Linking the nominals under this subtype is the high-telicity verb \textit{de} ‘reach.’ 

 \ea \label{ex:schaefer:sheg1}
\begin{xlist}
\ex \label{ex:schaefer:sheg1a}
\gll Òjè dé ókìn. \\
            Oje:\textsc{prx}	\textsc{pst}:reach:\textsc{pfv}	unconsciousness \\
           \glt  ‘Oje is unconscious.’
           
     \ex \label{ex:schaefer:sheg1b}
     \gll Òjè dé ɛ́khìì. \\
        Oje:\textsc{prx}   \textsc{pst}:reach:\textsc{pfv} crippledness \\
        \glt ‘Oje is crippled.’ 

    \ex \label{ex:schaefer:sheg1c}
     \gll Ójé dé ɛ̀hò. \\
        Oje:\textsc{prx}   \textsc{pst}:reach:\textsc{pfv} laziness \\
        \glt ‘Oje is lazy.’
    
\end{xlist}
\z

A second subtype differs from the \textsc{mda} based on non-distinctness of the entities involved in the predication. NP-1 and NP-2 do not refer to distinct entities. In the examples of (\ref{ex:schaefer:sheg2}) and (\ref{ex:schaefer:sheg3}), NP-1 is a human noun. NP-2 is a body part of NP-1. As a consequence, NP-1 and NP-2 do not stand as entities that are referentially distinct. Nonetheless, the examples in (\ref{ex:schaefer:sheg2}) and (\ref{ex:schaefer:sheg3}) differ in their meaning. In (\ref{ex:schaefer:sheg2}), for example, the verbs that link participants express a change of position: ‘pull,’ ‘reposition,’ or ‘extend.’ The NP-2 direct object for each is a body-part noun, respectively, \textit{éhɔ̀n} ‘ear,’ \textit{ɛ̀ò} ‘eye,’ and \textit{òɛ̀} ‘leg.’ 

 \ea \label{ex:schaefer:sheg2}
\begin{xlist}
\ex \label{ex:schaefer:sheg2a}
\gll Òjè yí éhɔ̀n. \\
            Oje:\textsc{prx}	\textsc{pst}:pull:\textsc{pfv}	ear \\
           \glt  ‘Oje is deaf.’

    \ex \label{ex:schaefer:sheg2b}
     \gll Òjè gbé ɛ̀ò. \\
        Oje:\textsc{prx}   \textsc{pst}:reposition:\textsc{pfv} eye \\
        \glt ‘Oje blinked.’

    \ex \label{ex:schaefer:sheg2c}
        \gll Ójé ɔ́ ɔ̀ fì òɛ̀.
        \\
        Oje:\textsc{dst}   \textsc{si:dst} \textsc{prs} extend:\textsc{ipfv} leg
        \\
        \glt ‘Oje limps.’
    
\end{xlist}
\z

The meanings assigned to the participants in the examples of (\ref{ex:schaefer:sheg3}) also do not refer to distinct entities. NP-2 is a body-part noun of the NP-1 participant. Although the meanings referred to involve a repositioning of a body part or the psychological state of a body, they all apply to a condition where the NP-1 referent is assumed to be interacting with the participants of a larger psycho-social/cultural matrix. Relative to this larger context of participants, the NP-1 referent is nosy, listening, or being patient.  

 \ea \label{ex:schaefer:sheg3}
\begin{xlist}
\ex \label{ex:schaefer:sheg3a}
\gll Ójé ɔ́ ɔ̀ hòò ɛ̀ò. \\
            Oje:\textsc{dst} \textsc{si:dst} \textsc{prs}	search:\textsc{ipfv}	eyes \\
           \glt  ‘Oje is nosy.’

    \ex \label{ex:schaefer:sheg3b}
     \gll Òjè gbé éhɔ́n kpɛ́kpɛ́kpɛ́. \\
        Oje:\textsc{prx}   \textsc{pst}:reposition:\textsc{pfv} ear attentively \\
        \glt ‘Oje listened attentively.’

    \ex \label{ex:schaefer:sheg3c}
        \gll Òjè zíɛ́n égbè.
        \\
        Oje:\textsc{prx}   \textsc{pst}:endure-\textsc{pfv} body
        \\
        \glt ‘Oje was patient.’
    
\end{xlist}
\z										

\textsc{tpp}s of a final subtype diverge from the \textsc{mda} principle in that NP-1 is an abstract instigator and NP-2 an affected human entity. NP-2 experiences the abstract condition instigated by NP-1. Nominals are linked by verbs of contact with fictive effect. 

\ea \label{ex:schaefer:sheg4}
\begin{xlist}
\ex \label{ex:schaefer:sheg4a}
    \gll ɛ̀khɔ̀ì ó òjè. \\
            shame:\textsc{prx} \textsc{pst}:enter:\textsc{pfv} Oje \\
           \glt  ‘Oje is ashamed.’

    \ex \label{ex:schaefer:sheg4b}
     \gll Ófɛ̀n ɔ̀ ɔ́ nwù òjè. \\
        fear:\textsc{prx}   \textsc{si:prx} \textsc{prs} take.hold:\textsc{ipfv} Oje \\
        \glt ‘Oje is becoming afraid.’

    \ex \label{ex:schaefer:sheg4c}
        \gll Úììn ɔ̀ ɔ́ gbè òjè.
        \\
        fever:\textsc{prx}   \textsc{si:prx} \textsc{prs} catch:\textsc{ipfv} Oje
        \\
        \glt ‘Oje is feverish.’

    \ex \label{ex:schaefer:sheg4d}
        \gll Òhànmì ɔ̀ ɔ́ gbè òjè. 
        \\
        hunger:\textsc{prx}   \textsc{si:prx} \textsc{prs} catch:\textsc{ipfv} Oje
        \\
        \glt ‘Oje is hungry.’   
\end{xlist}
\z	

With these examples we hope to have illustrated one feature of transitivity prominence in the Edoid language Emai. Across four sets of predications, members deviate from the transitive prototype but are still coded as transitive. We refer to these predications as peripheral members of the category transitive. Can we be satisfied with this general conclusion? There may be some room for doubt. 

How uniform are the predications in (\ref{ex:schaefer:sheg1}) - (\ref{ex:schaefer:sheg4}) in their failure to attest the transitive prototype and so to conflict with some definitional feature of transitivity? Suppose we say that the predications in sets (\ref{ex:schaefer:sheg1}) - (\ref{ex:schaefer:sheg4}) conflict in some fashion with our two definitions of transitivity: \citet{Naess2007} and \citet{Haspelmath2015}. First there is the principle of Maximally Distinguished Arguments (\textsc{mda}) advanced by \citet{Naess2007}. It held that a canonical transitive clause exhibited two participants with maximally distinct meanings. NP-1 should be a volitional instigator; NP-2 should be a non-volitional entity showing material effect. 

Across the four predication sets from Emai, there is no material effect, i.e., material change of state. There is no volitional instigator as NP-1 in any of the examples found in sets (\ref{ex:schaefer:sheg1}) - (\ref{ex:schaefer:sheg4}). Regarding the strong association of volition with human entities, there is a human noun as psychologically affected entity, NP-2 in set (\ref{ex:schaefer:sheg1}). It is not a non-human entity and it is not volitional for the predication in question. It is repeated as (\ref{ex:schaefer:sheg5}). 

 \ea \label{ex:schaefer:sheg5}
    \gll Òhànmì ɔ̀ ɔ́ gbè òjè. \\
            hunger:\textsc{prx} \textsc{si:prx} \textsc{prs} catch:\textsc{ipfv} Oje \\
           \glt  ‘Oje is hungry.’

\z	

None of the predications in (\ref{ex:schaefer:sheg1}) - (\ref{ex:schaefer:sheg4}) expresses change of material state of the ‘break’ type, which Haspelmath viewed as central to the transitive prototype. Perhaps physical effect among these predications is more like ‘hit,’ where the NP-2 argument is affected but does not undergo a material change of state. Or perhaps we can allow that each set expresses a very general sense of change vis-à-vis some sense of cultural normalcy, as in the instance of reaching a crippled condition: (\ref{ex:schaefer:sheg6}) repeated from the first subtype in (\ref{ex:schaefer:sheg1}). But each of these possible interpretations takes us well beyond a straightforward change in material effect.

 \ea \label{ex:schaefer:sheg6}
    \gll Òjè dé ɛ́khìì. \\
            Oje:\textsc{prx} \textsc{pst}:reach:\textsc{pfv} crippledness \\
           \glt  ‘Oje is crippled.’

\z	

\section{Transitivity scaling of verb meanings}

At this point it seems useful to review Haspelmath’s approach to the scaling of verb meanings by \citet{Tsunoda1985} and \citet{Malchukov2005}. Tsunoda provided a unidimensional scale of verb meanings and their likelihood of being coded as transitive. His scale can be represented as in \tabref{tab:schaefer:sheg7} with small letters representing verb meaning types and small caps indicating specific verbs that instantiate these meaning types. According to this scale, transitive coding is more likely for verb meanings on the left than for verb meanings on the right.  

\begin{table}
\caption{Unidimensional scale by Tsunoda for transitivity coding of verb types}
\label{tab:schaefer:sheg7}
 \begin{tabularx}{\textwidth}{llllllllll}
  \lsptoprule
{direct effect} & $>$ & perception & $>$ & pursuit & $>$ & cognition & $>$ & emotion \\
{\textsc{break:hit}} && {\textsc{see:look}} && \textsc{search} && \textsc{know} && \textsc{like} \\ 
  \lspbottomrule
 \end{tabularx}
\end{table}

Subsequently, Tsunoda’s unidimensional scale was restructured as predominantly two-dimensional by Malchukov, as in \tabref{tab:schaefer:sheg8}. He devised a split-scale for verb meanings that followed the meaning represented by \textsc{break}, which a host of prior studies associated with the transitive prototype \citep{Rice1987, Croft1993, Kittila2008}. Malchukov’s split-scale had an activity dimension and an experiential dimension, each of which consisted of different verb meaning types and the likelihood of their meaning being encoded as transitive. On the activity dimension, \textsc{hit} meanings were presented as more likely to be coded as transitive than \textsc{go} meanings. Between these two verb-meaning types, Malchukov positioned \textsc{search} meanings. Similarly, on the experiential dimension, the meanings for the \textsc{see} and \textsc{look} type were more likely to be coded as transitive than those for the \textsc{ache} type, with \textsc{fear} type meanings assigned a midpoint on the experiential scale. 

\begin{table}
\caption{Two-dimensional scale by Malchukov for transitivity coding of verb types}
\label{tab:schaefer:sheg8}
 \begin{tabularx}{.75\textwidth}{llllll}
  \lsptoprule
            & \textsc{hit} & \textsc{search}  & \textsc{go} & {(activity)}\\
  \textsc{break}  &      &       &           &  \\
     &   {\textsc{see:know}} &   \textsc{fear}  &    \textsc{ache}     & {(experiential)} \\
  \lspbottomrule
 \end{tabularx}
\end{table}

 Guided by Malchukov’s two-dimensional scale, Haspelmath computed the number of languages in the ValPaL database that coded a given verb meaning as transitive. The resulting numbers were then converted into percentage figures for each meaning type. The combination of meaning type and percentage as likelihood for transitivity coding were then arranged in the multi-layered scale shown in \tabref{tab:schaefer:sheg9}.  
 
\begin{table}
\caption{Multi-layered scale by Haspelmath for transitivity coding of verb types}
\label{tab:schaefer:sheg9}
 \begin{tabularx}{.75\textwidth}{llllll}
  \lsptoprule
            & {\textsc{hit} .94} & {\textsc{search} .89}  & {\textsc{go} .06}\\
  {\textsc{break} 1.0}  &      &       &           &  \\
     &   {\textsc{see} .92 : \textsc{know} .88} &   {\textsc{fear} .55}  &    {\textsc{ache} .12}\\
  \lspbottomrule
 \end{tabularx}
\end{table}

The verb-meaning ‘break’ in the ValPaL database was coded as transitive by 100\% of the thirty-five languages surveyed by Haspelmath. Activity-verb meanings aligned with \textsc{hit} and \textsc{search} were coded as transitive at the relatively high rates of 94\% and 89\%, respectively. On the other hand, the activity-verb meaning \textsc{go} in ValPaL was coded as transitive by only 6\% of the languages. On the experiential layer, we find that verb meanings related to \textsc{see} and \textsc{know}, respectively, were coded as transitive by 92\% and 88\% of the ValPaL languages. Verb meanings of the \textsc{fear} type manifested transitive coding in 55\% of the languages. Only 12\% coded verb meanings related to ‘ache’ as transitive. 

With respect to Haspelmath’s multi-layered scale, where do the Emai meanings expressed in example sets (\ref{ex:schaefer:sheg1}) - (\ref{ex:schaefer:sheg4}) fall? The predications in sets (\ref{ex:schaefer:sheg1}) - (\ref{ex:schaefer:sheg4}) fall exclusively onto the experiential dimension. Most reflect the extreme rightward point on the experiential dimension designated by \textsc{ache}, which we assume refers to experiential conditions that are internal to a human individual and pertain to that individual’s mental makeup. These conditions may have physical as well as psychological manifestations. This seems true for the predications in (\ref{ex:schaefer:sheg1}) where we find the meanings ‘be unconscious,’ ‘be crippled,’ and ‘be lazy.’ The predications in (\ref{ex:schaefer:sheg2}), ‘be deaf,’ ‘blink,’ and ‘limp,’ also appear to be experiential conditions associated with \textsc{ache}, although ‘blink’ seems not to have the negative or non-normative qualities associated with the other conditions in (\ref{ex:schaefer:sheg2}). As for the predications in (\ref{ex:schaefer:sheg3}), they align with the most leftward point on the experiential dimension. The verb meanings ‘be nosy,’ ‘listen,’ and ‘be patient,’ appear interactional and assume the involvement of other members or features of a culture. They express psycho-social conditions that operate with respect to other members of a society. The experiential scalar point most nearly aligned with these conditions is designated by \textsc{see} and \textsc{know}. As for the predications in (\ref{ex:schaefer:sheg4}), ‘be ashamed,’ ‘be afraid,’ ‘be hungry,’ and ‘be feverish,’ they align with the midpoint of the experiential dimension designated by \textsc{fear}. They represent conditions that are internal and most often non-normative relative to basic human experiences, but they are also particular to an individual’s mental makeup and explicitly call attention to their effect on an individual human.   

Much more could probably be said about each of these example sets and our initial attempt to distinguish them as subtypes. Experiential meanings of this sort are not often discussed in the semantically oriented literature of linguistics nor are the features or particulars that distinguish one experiential meaning from another given much attention, although meanings related to perception, most recently \citet{NorcliffMajid2024}, and emotion \citep{Lindquist2021} point in promising directions. Over the years much more linguistic analysis has been directed at verbs that capture events in the physical world rather than in the non-physical world. For the moment, we ask that the reader indulge us on this matter of analytic shortcomings. 

 Having tentatively assigned our Emai predication sets to various meaning points provided by Haspelmath’s experiential scale, we ask another question. Is the semantic space favorable to the transitivity coding of  meanings at the periphery limited to the experiential dimension?  

 What of the most rightward portion of the activity dimension, the scalar point designated by the verb form \textsc{go}? It, too, manifests a very low likelihood of transitivity coding since it was coded as transitive by only 6\% of the languages in ValPaL. 

Thus far we have not provided any evidence that transitivity coding in Emai is associated with the activity dimension, in particular activities related to meanings designated by \textsc{go}. We take this point on Haspelmath’s multi-layered scale to bear on activities of movement and motion where one entity moves through space, regardless of whether explicit mention is made of a second entity relative to which movement took place.  

To consider further the scalar point identified by \textsc{go} we turn to the likelihood of transitive coding for human body activities. A relevant domain concerns bodily emissions, which manifest two essential subtypes. Emissions from the human body can be characterized as either fluid/gaseous or sonorous (sound). In Emai, activities of bodily emission tend to be coded as transitive.  

The first activity type we consider denotes fluid or gaseous emissions. Emai meanings represented by such verbs as \textit{fɛna} ‘excrete,’ \textit{roo} ‘release,’ \textit{vbia} ‘discharge,’ \textit{nɛ} ‘pass,’ and \textit{fi} ‘exhale’ articulate as transitive structures \REF{ex:schaefer:sheg10}. They combine with nouns in direct object position, respectively \textit{ìsɔ̀n} ‘feces,’ \textit{áàhìɛ̀n} ‘urine,’ \textit{évìɛ̀} ‘tears,’ \textit{èsɛ̀ìn} ‘spittle,’ \textit{ìhɔ̀n} ‘fart,’ and \textit{étìn} ‘breadth.’ Each verb-noun combination in construction with a subject noun refers to the emission of a fluid or gaseous substance and encodes as transitive. 

 \ea \label{ex:schaefer:sheg10}
\begin{xlist}
\ex \label{ex:schaefer:sheg10a}
    \gll Òjè fɛ́ná ìsɔ̀n. \\
          Oje:\textsc{prx} \textsc{pst}:excrete:\textsc{pfv} feces \\
           \glt  ‘Oje has excreted his feces/has defecated.’

    \ex \label{ex:schaefer:sheg10b}
    \gll Òjè fɛ́ná áàhìɛ̀n. \\
          Oje:\textsc{prx}   \textsc{pst}:excrete:\textsc{pfv} urine \\
           \glt  ‘Oje has excreted his urine/has urinated.’

    \ex \label{ex:schaefer:sheg10c}
    \gll Àlèkè róó évìɛ̀. \\
          {Aleke:\textsc{prx}} {\textsc{pst}:release:\textsc{pfv}} tears
        \\
           \glt  ‘Aleke has shed tears/has teared up.’

    \ex \label{ex:schaefer:sheg10d}
    \gll Òjè vbíá èsɛ̀ìn.  \\
          Oje:\textsc{prx}   {\textsc{pst}:discharge:\textsc{pfv}} spittle
        \\
           \glt  ‘Oje has discharged spittle/has spat.’

    \ex \label{ex:schaefer:sheg10e}
    \gll Òjè nɛ́ ìhɔ̀n.  \\
          Oje:\textsc{prx}   {\textsc{pst}:pass:\textsc{pfv}} fart       \\
           \glt ‘Oje passed a fart/has farted.’

    \ex \label{ex:schaefer:sheg10f}
    \gll Òjè ɔ̀ ɔ́ fì étìn.  \\
         Oje:\textsc{prx}   {\textsc{si:prx}} \textsc{prs} {exhale:\textsc{ipfv}} breath      \\
           \glt ‘Oje is breathing.’
    
\end{xlist}
\z	

Second, we identify expressions for activities that code a sound emission. In Emai some of the relevant verbs include \textit{ta} ‘speak,’ \textit{so} ‘sing,’ and \textit{zɛ} ‘express.’ Each occurs with a following noun in direct object position, sometimes with a cognate noun (\textit{étà} ‘words’) but not usually (\textit{íòò} ‘song,’ or ‘thought,’ \textit{ùnyɔ̀} ‘grumble’). Together with a subject noun, these verb-noun combinations are coded in a transitive pattern.  

 \ea \label{ex:schaefer:sheg11}
\begin{xlist}
\ex \label{ex:schaefer:sheg11a}
    \gll Òjè tá étà. \\
        Oje:\textsc{prx} \textsc{pst}:speak:\textsc{pfv} words \\
           \glt ‘Oje has spoken.’

\ex \label{ex:schaefer:sheg11b}
    \gll Òjè ɔ̀ ɔ́ sò íòò. \\
        Oje:\textsc{prx}   \textsc{si:prx} \textsc{prs} {sing:\textsc{ipfv}} song \\
           \glt ‘Oje is singing a song/ singing.’

\ex \label{ex:schaefer:sheg11c}
    \gll Òjè ɔ̀ ɔ́ zɛ̀ ùnyɔ̀. \\
        {Oje:\textsc{prx}} \textsc{si:prx} \textsc{prs} {express:\textsc{ipfv}} grumble \\
           \glt ‘Oje is grumbling.’

\ex \label{ex:schaefer:sheg11d}
    \gll Yàn á zɛ̀ ìòò.  \\
        {\textsc{3pl:prx}} \textsc{prs} {express:\textsc{ipfv}} thought \\
           \glt ‘They are conversing.’
    
\end{xlist}
\z	

To be sure there are bodily emission meanings that are coded as intransitive. A few of the relevant verbs are seen in (\ref{ex:schaefer:sheg12}), where we find the meanings for \textit{tiho} ‘sneeze’ and \textit{oo} ‘ooze’ coded as intransitive. They do not participate in transitive coding patterns.  

 \ea \label{ex:schaefer:sheg12}
\begin{xlist}
\ex \label{ex:schaefer:sheg12a}
    \gll Òjè tíhó-ì. \\
        Oje:\textsc{prx} {\textsc{pst}:sneeze-\textsc{pfv}} \\
           \glt ‘Oje has sneezed.’

\ex \label{ex:schaefer:sheg12b}
    \gll Élì ɛ̀màì óó-ì. \\
        \textsc{art} {boil:\textsc{prx}} {\textsc{pst}:ooze-\textsc{pfv}} \\
           \glt ‘The boils have oozed.’
    
\end{xlist}
\z	

Transitive coding relevant to Haspelmath’s activity dimension is not limited to expressions of bodily emission. There are motion activities, specifically of the motion+path type, that Emai codes as transitive. Representative examples are shown in (\ref{ex:schaefer:sheg13}). They include the motion+path verbs \textit{ye} ‘move toward,’ \textit{lagaa} ‘move around, encircle,’ \textit{shan} ‘move through,’ \textit{hɛɛn} ‘ascend,’ and \textit{kpɔɔn} ‘descend.’ Each combines with a noun of location in direct object position as well as a subject noun.  

 \ea \label{ex:schaefer:sheg13}
\begin{xlist}
\ex \label{ex:schaefer:sheg13a}
    \gll Àlèkè yé èkó. \\
        {Aleke:\textsc{prx}} {\textsc{pst}:move.toward:\textsc{pfv}} Lagos \\
           \glt ‘The woman has moved toward Lagos.’

\ex \label{ex:schaefer:sheg13b}
    \gll Àlèkè ɔ̀ ɔ́ làgáá ùhàì. \\
        {Aleke:\textsc{prx}} {\textsc{si:prx}} \textsc{prs} {move.around:\textsc{ipfv}} well \\
           \glt ‘Aleke is moving around / circling the well.’

\ex \label{ex:schaefer:sheg13c}
    \gll Àlèkè shán égbóà. \\
        {Aleke:\textsc{prx}} {\textsc{pst}:move.through:\textsc{pfv}} backyard \\
           \glt ‘Aleke has moved through the backyard.’

\ex \label{ex:schaefer:sheg13d}
    \gll Àlèkè hɛ́ɛ́n ɔ́lí ɔ́kɔ̀ɔ́. \\
        {Aleke:\textsc{prx}} {\textsc{pst}:ascend:\textsc{pfv}} \textsc{art} hill \\
           \glt ‘Aleke has ascended the hill.’

\ex \label{ex:schaefer:sheg13e}
    \gll Àlèkè kpɔ́ɔ́n ɔ́lí ɔ́kɔ̀ɔ́. \\
        {Aleke:\textsc{prx}} {\textsc{pst}:descend:\textsc{pfv}} \textsc{art} hill \\
           \glt ‘Aleke has descended the hill.’

\end{xlist}
\z	

The transitive coding of the motion+path verbs in (\ref{ex:schaefer:sheg13}) contrasts with the intransitive coding shown by the Emai verb corresponding to the designator point \textsc{go} on Haspelmath’s activity dimension. Corresponding to \textsc{go} is the Emai verb \textit{lodɛ} ‘go’ that is coded as intransitive. Its following noun phrase is coded as oblique by locative marker \textit{vbi}.  

\ea \label{ex:schaefer:sheg14}
    \gll Àlèkè lòdɛ́ vbì ìwè. \\
        Aleke:\textsc{prx} \textsc{prs}:go:\textsc{ipfv} \textsc{loc} house\\
           \glt ‘Aleke is going to the house.’
\z	

Emai shows a similar coding pattern that employs oblique noun phrases for a few motion+path verbs. They are presented in (\ref{ex:schaefer:sheg15}). Both \textit{o} ‘enter’ and \textit{sɛ} ‘move as far as’ are coded as intransitive, as was the Emai counterpart of \textsc{go}.  

 \ea \label{ex:schaefer:sheg15}
\begin{xlist}
\ex \label{ex:schaefer:sheg15a}
    \gll Àlèkè ó vbí ɛ́kɔ́à. \\
        {Aleke:\textsc{prx}} {\textsc{pst}:enter:\textsc{pfv}} \textsc{loc} room \\
           \glt ‘Aleke has entered the room.’

\ex \label{ex:schaefer:sheg15b}
    \gll Àlèkè sɛ́ vbí ɛ́dà. \\
        {Aleke:\textsc{prx}} {\textsc{pst}:move.up.to:\textsc{pfv}} \textsc{loc} river \\
           \glt ‘Aleke has moved as far as the river.’

\end{xlist}
\z	

\section{Discussion}

Given the preceding, one can ask what definitional elements remain from our attempt to identify predications at the periphery of transitivity in Emai? Recall that these elements were framed by the principle of Maximally Distinguished Arguments (\textsc{mda}), which held that a canonical transitive clause exhibited two participants with maximally distinct meanings. NP-1 was a volitional instigator; NP-2 was a non-volitional, affected entity displaying material effect. Neither of these conditions on the status of transitive clause participants held for the simplex predications of Emai bodily-conditions reviewed in (\ref{ex:schaefer:sheg1}) - (\ref{ex:schaefer:sheg4}). It appears that it is only the general condition on form that remains of the \textsc{mda}: A transitive verb must code a predication with two participants. 

Let’s consider another possibility. Could the examples or perhaps some of the examples in sets (\ref{ex:schaefer:sheg1}) - (\ref{ex:schaefer:sheg4}) belong to a grammatical domain other than transitivity? The copula domain comes to mind. Although copula status is perhaps far-fetched, the English translations of most examples in (\ref{ex:schaefer:sheg1}) - (\ref{ex:schaefer:sheg4}) employ a form of BE. 

 Emai has five constructions in its copula domain. Respectively, they assign a nominal to a state of existence (\ref{ex:schaefer:sheg16a}), clarify the discourse relation between copula nominals (\ref{ex:schaefer:sheg16b}), assign a nominal to a property (\ref{ex:schaefer:sheg16c}), assign a nominal to a location (\ref{ex:schaefer:sheg16d}), and assign a nominal to a class (\ref{ex:schaefer:sheg16e}).  

  \ea \label{ex:schaefer:sheg16}
\begin{xlist}
\ex \label{ex:schaefer:sheg16a}
    \gll Óì kí ɔ́ɔ̀. \\
        thief \textsc{ncf} be\\
           \glt ‘It isn't a thief.’

\ex \label{ex:schaefer:sheg16b}
    \gll ɔ́áìn lí í khì àgbògbòràn. \\
        that.one \textsc{pcf} \textsc{id} be woodpecker \\
           \glt ‘That one is a woodpecker.’

\ex \label{ex:schaefer:sheg16c}
    \gll ɔ́lì ɛ̀kpɛ̀n ú nwɛ̀nɛ́nwɛ̀nɛ́. \\
        \textsc{art} {leopard:\textsc{prx}} {\textsc{pst}:be:\textsc{pfv}} spotted \\
           \glt ‘The leopard is spotted.’

\ex \label{ex:schaefer:sheg16d}
    \gll ɔ́lí ɔ́vbèkhàn ríì vbí ìwè. \\
        \textsc{art} {youth:\textsc{prx}} {\textsc{pst}.be} \textsc{loc} house \\
           \glt ‘The youth is in the house.’

\ex \label{ex:schaefer:sheg16e}
    \gll ɔ́lí ɔ́mɔ̀hè í ì vbì óì. \\
        \textsc{art} man \textsc{si} \textsc{neg} be thief \\
           \glt ‘The man is not a thief.’
           
\end{xlist}
\z	

Not all \textsc{be} constructions provide an acceptable syntactic frame for the verbs in sets (\ref{ex:schaefer:sheg1}) - (\ref{ex:schaefer:sheg4}) or the sets that came later in the paper. Four of the five \textsc{be} forms exhibit a two-participant structure. The specifier has only one participant (\ref{ex:schaefer:sheg16a}). One of the five with a two-participant structure requires a pronoun of identity as participant one and relates this pronoun to a preceding nominal in focus position. This complex structure is obligatory for the equational identity copula (\ref{ex:schaefer:sheg16b}). Two of the five limit the second participant to a particular type of semantic category: property term for the adjective in (\ref{ex:schaefer:sheg16c}) and locative for the spatial referent in (\ref{ex:schaefer:sheg16d}). One copula, class membership in (\ref{ex:schaefer:sheg16e}), shows two-participant nominal structures. However, it exhibits a strong preference for negative polarity and displays affirmative structures only in exceptional discourse contexts.  

No copula seems to fit. None of the restrictions on copulas hold for the predication sets (\ref{ex:schaefer:sheg1}) - (\ref{ex:schaefer:sheg4}) from Emai or, for that matter, the predication sets in (\ref{ex:schaefer:sheg5}) through (\ref{ex:schaefer:sheg15}). 

Copula assignment might be considered for set (\ref{ex:schaefer:sheg1}), which involved verb \textit{de} ‘reach.’ But \textit{de} is not a copula. Normatively, it allows two-participant structures, where the second participant is a location, as in (\ref{ex:schaefer:sheg17}).  

  \ea \label{ex:schaefer:sheg17}
    \gll Òjè dé ókhúnmí ɔ́kɔ̀ɔ́. \\
         Oje:\textsc{prx} \textsc{pst}:reach:\textsc{pfv} top hill \\
           \glt ‘Oje reached the top of the hill.’
\z	
 
Perhaps one could propose that \textit{de} is a copula in the making. But copular status is not where the meaning of \textit{de} stands today. It seems time to conclude that there is no other grammatical domain other than the periphery of transitivity to which sets (\ref{ex:schaefer:sheg1}) - (\ref{ex:schaefer:sheg4}) might belong.   

Moving on, the analysis undertaken in this paper, although leaving some matters unresolved, directs attention to potential areas of future investigation concerning transitivity coding. What we know from the present analysis of Emai is that the patterns of transitivity coding are not limited to one side of the scalar dimensions presented by Malchukov and Haspelmath. That is the side closest to \textsc{break}. 

Furthermore, we have seen that Emai transitivity coding is found across both scalar dimensions: activity and experiential. Relatively often in the past it has been assumed that transitivity coding is less likely for items in the experiential domain where meanings for mental events can be identified \citep{Croft1993}. Nonetheless, the way transitivity coding for mental experiences diverges from the \textsc{break} prototype has not been explicitly addressed in the linguistic literature. Nor has the literature explored the manner of divergence for non-mental experiences, i.e. physical activities such as motion+path verbs. Both strategies of divergence point to potential areas of future investigation. 

For the moment let us return to the \textsc{break} prototype. It consists of the elements proposed by the \textsc{mda} Principle of Næss: two participants, participants that are maximally distinct (volitional instigator as NP-1 and non-volitional affected entity as NP-2), and realization of a material effect or material change of state.  

Which of these elements is more susceptible to failure for verb meanings that do not reflect the transitivity prototype. In the case of Emai these meanings are on the experiential dimension or the activity dimension. Do all transitivity elements fail uniformly and across the board? Might there be a certain asymmetry among these elements of transitivity in their failed application to the activity dimension as opposed to the experiential dimension? Answers to such questions may provide us with additional insight into the nature of linguistic transitivity and for that matter intransitivity. 

Generalizations that can be drawn from the present study of Emai are two-fold. They pertain to the features of maximal distinctness and realization of material effect. Experiential predications failed regarding the \textsc{mda} features maximal distinctness and realization of material effect. Both were absent. Activity predications failed concerning the \textsc{mda} feature realization of material effect. Regarding the feature maximal distinctness, activity predications showed a split. Meanings for motion+path verbs manifested maximally distinct participants, whereas meanings for bodily emission verbs did not. Motion+path verbs (e.g. \textit{hɛɛn} ‘ascend,’ \textit{lagaa} ‘encircle’) that were coded as transitive exhibited a moving object as NP-1 and a location as NP-2. The referents for these noun phrases are arguably maximally distinct, especially when compared to bodily emission verbs. The latter, for example \textit{fɛna} ‘excrete,’ and \textit{vbia} ‘discharge,’ employ noun phrases that are not maximally distinct. NP-1 refers to a human noun and NP-2 refers to a substance that is retained within that human referent. In this respect, bodily emission verbs are comparable to bodily condition verbs, e.g. \textit{fì òɛ̀} ‘limp’ from ‘extend leg.’  

Additionally, the realization of material effect feature relates in an interesting way to activity verb meanings. Neither activity type, motion+path or bodily emission, denotes a material change of state, a material effect. However, each refers to a change of location for an entity: NP-1 for verbs of motion+path and NP-2 for verbs of bodily emission. Furthermore, NP-1 for motion+path and bodily emission verbs could be viewed as a volitional instigator, or at least instigator, of their respective changes of location. Common to both the transitive prototype and the verbs of motion+path activity is change. Change of state in the former and change of location in the latter. Viewed from this perspective, the propensity for transitivity coding on the activity dimension seems less distant from the \textsc{mda} prototype than does transitivity coding on the experiential dimension. 

The preceding discussion has focused exclusively on transitive coding of meanings. Intransitive coding has remained in the background. A strong focus on intransitivity, however, has the potential to increase our understanding of transitivity. For example, little sustained linguistic investigation has been directed toward whether there are meanings confined exclusively to intransitive coding within and across languages. In other words, are there bedrock meanings that code as intransitive, just as there are bedrock meanings that code as transitive, e.g. \textsc{break}? It is our belief that meanings that cross-linguistically tend to show exclusively intransitive coding might provide a greater understanding of how feature complexes underlying transitivity, such as those designed for the \textsc{break} type, may be extended beyond the transitive prototype and so to the way transitivity is characterized.  

As a final comment, the Emai data sets suggest an additional line of inquiry. Further investigation of transitivity prominence in West Africa, especially for meanings aligned with the experiential dimension, might usefully question whether nominals other than those for body parts or abstractions contribute to the failure of the asymmetric force-dynamic relationship between participants that is the linchpin of the transitive prototype \citep{Rice1987}.  

\section*{Abbreviations}
\begin{tabularx}{.5\textwidth}{@{}lQ@{}}
\textsc{art} & article\\
\textsc{id} & identity pronoun\\
\end{tabularx}%
\begin{tabularx}{.5\textwidth}{@{}lQ@{}}
\textsc{dst} & distal temporal distance\\
\textsc{impf} & imperfective viewpoint aspect\\
\end{tabularx}
\begin{tabularx}{.5\textwidth}{@{}lQ@{}}
\textsc{loc} & locative\\
\textsc{ncf} & negative contrastive focus\\
\end{tabularx}
\begin{tabularx}{.5\textwidth}{@{}lQ@{}}
\textsc{neg} & negation\\
\textsc{pcf} & positive contrastive focus\\
\end{tabularx}
\begin{tabularx}{.5\textwidth}{@{}lQ@{}}
\textsc{pfv} & perfective viewpoint aspect\\
\textsc{prs} & present tense\\
\end{tabularx}
\begin{tabularx}{.5\textwidth}{@{}lQ@{}}
\textsc{prx} & proximal temporal distance\\
\textsc{pst} & past tense\\
\end{tabularx}
\begin{tabularx}{.5\textwidth}{@{}lQ@{}}
\textsc{si} & subject index\\
 \\
\end{tabularx}

\section*{Acknowledgements}
Data incorporated in this paper is derived from research fieldwork over several decades that involved narrative texts, dictionary, and grammar construction. This research was sponsored by the National Science Foundation, (BNS \#9011338 and SBR \#9409552), U.S. Department of State (CUAP ASJY 1333), and the National Endowment for the Humanities (PD-50004-06). 

\sloppy
\printbibliography[heading=subbibliography,notkeyword=this]
\end{document}
